\section*{Hoofdstuk \ref{ch:b1:field-lab-methods-1} --- Veldwerk- en laboratoriummethoden I}

\begin{solution}{exo:b1:field-lab-methods-1:1}
Op 1:10\,000 wordt \qty{4.9}{m} $4.9/10\,000 = \qty{0.49}{mm}$; op
1:25\,000, $4.9/25\,000 = \qty{0.20}{mm}$, ongeveer de breedte van een potloodlijn.
\end{solution}

\begin{solution}{exo:b1:field-lab-methods-1:2}
In Reykjavik is $D = \ang{-11.28}$: $52 - 11.28 = \ang{040.7}$. In Wellington is
$D = \ang{+23.23}$: $52 + 23.23 = \ang{075.2}$. Eenzelfde kompasaflezing betekent \omterm{def:b1:field-lab-methods-1:trend-plunge}{richtingen} die
\ang{34.5} uit elkaar liggen.
\end{solution}

\begin{solution}{exo:b1:field-lab-methods-1:3}
$20\to015$. Met $t = \ang{15}$ en $p = \ang{20}$ is de neerwaartse eenheidsvector
\[
  (\cos\ang{20}\sin\ang{15},\ \cos\ang{20}\cos\ang{15},\ -\sin\ang{20}) =
  (0.2432,\ 0.9077,\ -0.3420).
\]
\end{solution}

\begin{solution}{exo:b1:field-lab-methods-1:4}
Het rekenkundig gemiddelde is $(350 + 358 + 6)/3 = \ang{238}$, bijna tegengesteld aan alle drie
aflezingen. De horizontale eenheidsvectoren $(\sin a, \cos a)$ tellen op tot
$(-0.1736 - 0.0349 + 0.1045,\ 0.9848 + 0.9994 + 0.9945) = (-0.1040, 2.9787)$,
waarvan het \omterm{def:b1:spherical-astronomy:horizontal}{azimut} $\operatorname{atan2}(-0.1040, 2.9787) = \ang{-2.0}$ is, dat is
\ang{358.0}: het juiste gemiddelde, in het midden van de aflezingen.
\end{solution}

\begin{solution}{exo:b1:field-lab-methods-1:5}
Gemiddelde $175/5 = \ang{35.0}$. Afwijkingen $-1, 2, 0, -2, 1$; som van de kwadraten 10;
$s = \sqrt{10/4} = \ang{1.58}$. Standaardfout $1.58/\sqrt5 = \ang{0.71}$.
\end{solution}

\begin{solution}{exo:b1:field-lab-methods-1:6}
De normalen zijn $(\sin\ang{30}, 0, \cos\ang{30}) = (0.5, 0, 0.866)$ en
$(-0.5, 0, 0.866)$; $\vec S = (0, 0, 1.732)$ is verticaal: het gemiddelde is een horizontaal vlak, en
$\bar R = 1.732/2 = 0.866$, elke aflezing \ang{30} van het gemiddelde. De twee aflezingen bemonsteren niet
één oppervlak met fouten: zij behoren tot twee poten van een \omterm{def:b1:geological-maps:fold}{plooi}, twee verschillende populaties, en hun
gemiddelde beschrijft geen van beide. Elke poot wordt afzonderlijk gemiddeld; samen bepalen zij de \omterm{def:b1:geological-maps:fold}{plooias},
de lijn die beide vlakken gemeen hebben, hier horizontaal en noord--zuid.
\end{solution}

\begin{solution}{exo:b1:field-lab-methods-1:7}
$\tan p = \tan\ang{12}\cos(\ang{80} - \ang{120}) = 0.2126\times0.7660 =
0.1628$: $p = \ang{9.2}$, de krassen \omterm{def:b1:field-lab-methods-1:trend-plunge}{duiken} naar \ang{120}. Voor de \omterm{def:b1:field-lab-methods-1:trend-plunge}{richting} \ang{170} is
$\cos(\ang{-90}) = 0$: de lijn is horizontaal, langs de \omterm{def:b1:geological-maps:attitude}{strijk}.
\end{solution}

\begin{solution}{exo:b1:field-lab-methods-1:8}
$0.47/0.20 = 2.35$ jaar na 1 januari 2026, dat is in het voorjaar van 2028. Op
1 januari 2030: $-0.47 + 4\times0.20 = \ang{+0.33}$, oostelijk van het geografische noorden.
\end{solution}

\begin{solution}{exo:b1:field-lab-methods-1:9}
De drie normalen liggen in het verticale vlak met \omterm{def:b1:spherical-astronomy:horizontal}{azimut} \ang{080}, op \ang{10}, \ang{12} en \ang{14} van de
verticale; de buitenste twee zijn symmetrisch om de middelste, dus hun som ligt er langs: het gemiddelde is
$080/12$, met $\theta = \ang{2}, 0, \ang{2}$. Dan is $\bar R = (1 +
2\cos\ang{2})/3 = (1 + 1.998782)/3 = 0.999594$, en $\sqrt{2(1 - \bar R)} =
\sqrt{0.000812} = \qty{0.02850}{rad} = \ang{1.633}$; exact,
$\theta_{\text{rms}} = \sqrt{(4 + 0 + 4)/3} = \ang{1.633}$.
\end{solution}

\begin{solution}{exo:b1:field-lab-methods-1:10}
De doorsnede is een cirkel met straal $\sqrt{r^2 - h^2}$. Haar gemiddelde is
$\frac1r\int_0^r\sqrt{r^2 - h^2}\,\mathrm dh = \frac1r\cdot\frac{\pi r^2}{4}
= \frac{\pi r}{4}$, omdat de integraal de oppervlakte van een kwartcirkel is. Korrels van
\qty{0.50}{mm} in doorsnede lijken gemiddeld $0.50\times\pi/4 = \qty{0.39}{mm}$ in doorsnede:
een \omterm{def:b1:field-lab-methods-1:thin-section}{dunne slijpplaat} onderschat de korrelgrootte.
\end{solution}

\begin{solution}{exo:b1:field-lab-methods-1:11}
De draaiing over $D$ om de verticale,
$(x, y, z) \mapsto (x\cos D + y\sin D,\ -x\sin D + y\cos D,\ z)$, stuurt
$(\sin\delta\sin\alpha, \sin\delta\cos\alpha, \cos\delta)$ naar
$(\sin\delta\sin(\alpha + D), \sin\delta\cos(\alpha + D), \cos\delta)$: zij telt $D$ bij elke
\omterm{def:b1:geological-maps:attitude}{hellingrichting}. Omdat zij lineair is, stuurt zij de resultante $\vec S$ naar de resultante van de gedraaide normalen;
zij behoudt lengten en de verticale component, dus de gemiddelde helling $\arccos(S_U/\|\vec S\|)$ verandert niet, en zij
draait het horizontale deel van $\vec S$ over $D$, zodat de gemiddelde \omterm{def:b1:geological-maps:attitude}{hellingrichting} met $D$ toeneemt. Elke aflezing
corrigeren voor de declinatie, of alleen hun gemiddelde, geeft hetzelfde resultaat.
\end{solution}

\begin{solution}{exo:b1:field-lab-methods-1:12}
De spreiding langs de lijn van de steilste helling is een spreiding van de helling: de standaardfout is
$2/\sqrt N \le 0.5$ voor $N \ge 16$. Dwars erop verandert een draaiing van de normaal over \ang{2} de \omterm{def:b1:geological-maps:attitude}{hellingrichting} met
$2/\sin\ang{12} = \ang{9.62}$; de standaardfout van de gemiddelde \omterm{def:b1:geological-maps:attitude}{hellingrichting} is $9.62/\sqrt N \le 1$ voor
$N \ge 92.5$, dat is 93 aflezingen. Een flauwe helling is snel bekend; haar \omterm{def:b1:field-lab-methods-1:trend-plunge}{richting} niet.
\end{solution}

\begin{solution}{pb:b1:field-lab-methods-1:1}
\textbf{1.} \omterm{def:b1:geological-maps:attitude}{Hellingrichtingen}: $967/12 = \ang{80.6}$; hellingen: $145/12 =
\ang{12.1}$.

\textbf{2.} \ang{355} en \ang{005} middelen tot \ang{180}, de tegengestelde
\omterm{def:b1:field-lab-methods-1:trend-plunge}{richting}. Hier liggen alle \omterm{def:b1:geological-maps:attitude}{hellingrichtingen} tussen \ang{076} en \ang{085}, ver weg van de sprong bij
\ang{000}, dus het rekenkundig gemiddelde werkt.

\textbf{3.} $\|\vec n\|^2 = \sin^2\delta(\sin^2\alpha + \cos^2\alpha) +
\cos^2\delta = 1$. De neerwaartse hellingslijn heeft de \omterm{def:b1:field-lab-methods-1:trend-plunge}{richting}
$(\cos\delta\sin\alpha, \cos\delta\cos\alpha, -\sin\delta)$, en haar inproduct met $\vec n$ is
$\sin\delta\cos\delta - \cos\delta\sin\delta = 0$; de strijklijn is horizontaal en loodrecht op de
\omterm{def:b1:geological-maps:attitude}{hellingrichting}, $(\cos\alpha, -\sin\alpha, 0)$, en haar inproduct met $\vec n$ is
$\sin\delta\sin\alpha\cos\alpha - \sin\delta\cos\alpha\sin\alpha = 0$.

\textbf{4.} Aflezing 11 is dezelfde als aflezing 2: $(0.2222, 0.0352, 0.9744)$.
Aflezing 12:
\[
  (\sin\ang{11}\sin\ang{82},\ \sin\ang{11}\cos\ang{82},\ \cos\ang{11}) =
  (0.1890,\ 0.0266,\ 0.9816).
\]

\textbf{5.} $\vec S = (2.4747,\ 0.4082,\ 11.7294)$.

\textbf{6.} $\|\vec S\| = \sqrt{6.1241 + 0.1666 + 137.5788} = 11.99457$, en
$\bar R = 11.99457/12 = 0.999547$.

\textbf{7.} $\bar\alpha = \operatorname{atan2}(2.4747, 0.4082) = \ang{80.6}$
en $\bar\delta = \arccos(11.7294/11.99457) = \arccos 0.977893 = \ang{12.1}$:
het vectorgemiddelde is $080.6/12.1$, hetzelfde als de rekenkundige gemiddelden hier.

\textbf{8.} \omterm{def:b1:geological-maps:attitude}{Hellingrichtingen}: $s = \ang{2.94}$; hellingen: $s = \ang{1.68}$.

\textbf{9.} $\theta_{\text{rms}} \approx \sqrt{2\times0.000453} =
\qty{0.0301}{rad} = \ang{1.72}$.

\textbf{10.} $\sqrt{1.68^2 + (\sin\ang{12.1}\times2.94)^2} = \sqrt{2.82 +
0.38} = \ang{1.79}$, dicht bij \ang{1.72}: de aflezingen gedragen zich zoals het foutenmodel zegt, binnen de
effecten van de afronding op de graad en van een kleine steekproef.

\textbf{11.} De helling: zij draagt 2.82 van de 3.20 bij (in graden in het kwadraat). De \omterm{def:b1:geological-maps:attitude}{hellingrichting} telt
alleen mee via $\sin\delta = 0.21$, omdat een verandering van \omterm{def:b1:geological-maps:attitude}{hellingrichting} op een flauwe laag het vlak zeer weinig draait.

\textbf{12.} Gemiddelde helling: $1.68/\sqrt{12} = \ang{0.48}$; gemiddelde \omterm{def:b1:geological-maps:attitude}{hellingrichting}:
$2.94/\sqrt{12} = \ang{0.85}$.

\textbf{13.} Nee: \ang{2.9} \omterm{def:b1:geological-maps:attitude}{hellingrichting} op een laag van \ang{12} is slechts
$2.9\sin\ang{12.1} = \ang{0.6}$ draaiing van het vlak, minder dan de \ang{1.7} van de hellingen. Elke lichte golving
van een flauw gelaagdheidsvlak spreidt de \omterm{def:b1:geological-maps:attitude}{hellingrichting} ruim.

\textbf{14.} $D = \ang{-2.5}$: $83.1 - 2.5 = \ang{080.6}$; haar gemiddelde is
$080.6/12.1$, hetzelfde als dat van de eerste geoloog.

\textbf{15.} De helling wordt met de kliseometer afgelezen, tegen de verticale die door de zwaartekracht wordt
gegeven; de naald speelt geen rol.

\textbf{16.} Ja: een draaiing om de verticale verplaatst elke aflezing en het gemiddelde met dezelfde hoek
(\cref{exo:b1:field-lab-methods-1:11}); het antwoord is hetzelfde.

\textbf{17.} $10\times0.2 = \ang{2}$ fout op elke \omterm{def:b1:geological-maps:attitude}{hellingrichting}, meer dan twee keer de standaardfout van \ang{0.85}:
een systematische fout die het middelen niet kan opsporen.

\textbf{18.} $80.6 - 90 = \ang{350.6}$; de oriëntatie is \omterm{def:b1:geological-maps:attitude}{strijk} \ang{351}, helling \ang{12} naar het oosten.

\textbf{19.} \omterm{def:b1:geological-maps:attitude}{Hellingrichtingen} $080.6$ en $080.6$: geen verschil op deze precisie; hellingen $12.1$ en $12.2$: \ang{0.1}.

\textbf{20.} De normalen $(\sin\ang{12.2}\sin\ang{80.6},
\sin\ang{12.2}\cos\ang{80.6}, \cos\ang{12.2})$ en $\vec S/\|\vec S\|$ maken
een hoek van \ang{0.13}.

\textbf{21.} $1.72/\sqrt{12} = \ang{0.50}$; het verschil, \ang{0.13}, is een kwart daarvan: niet significant.
Het kompas en de boringen komen overeen.

\textbf{22.} Vier keer zoveel: 48 aflezingen.

\textbf{23.} Bijvoorbeeld: een verkeerde of verouderde declinatie-instelling; een kompas dat door ijzer in de
buurt wordt afgebogen (een voertuig, gereedschap, magnetische gesteenten); een kliseometer die niet meer
gekalibreerd is; aflezingen die alleen op de best ontsloten vlakken zijn gedaan, die de poot niet hoeven te
vertegenwoordigen. Middelen verkleint toevallige fouten, niet systematische.

\textbf{24.} De \omterm{def:b1:field-lab-methods-1:mean-orientation}{gemiddelde oriëntatie} van de oostflank is \textbf{081/12}: \omterm{def:b1:geological-maps:attitude}{hellingrichting}
\ang{081}, helling \ang{12}.
\end{solution}
